\documentclass[12pt, a4paper]{article}

\usepackage[utf8]{inputenc}
\usepackage[T1]{fontenc}
\usepackage[brazil]{babel}
\usepackage{amsmath, amssymb, amsthm}
\usepackage{geometry}
\usepackage{booktabs}
\usepackage{array}
\usepackage{xcolor}
\usepackage{mdframed}
\usepackage{hyperref}

\geometry{margin=2.5cm}

\definecolor{destaque}{RGB}{0, 90, 160}
\definecolor{cinza}{RGB}{245, 245, 245}

\newmdenv[
  backgroundcolor=cinza,
  linecolor=destaque,
  linewidth=1.5pt,
  roundcorner=4pt,
  innertopmargin=10pt,
  innerbottommargin=10pt
]{caixadestaque}

\title{
  \textbf{\textcolor{destaque}{O Problema do XOR}}\\[0.4em]
  \large Demonstração Completa via Equação Normal
}
\author{}
\date{}

\begin{document}

\maketitle
\thispagestyle{empty}

% ─────────────────────────────────────────────
\section*{1. Configuração do Problema}

O dataset XOR é composto por quatro amostras:

\begin{center}
\begin{tabular}{ccc}
\toprule
$x_1$ & $x_2$ & $y$ \\
\midrule
0 & 0 & 0 \\
0 & 1 & 1 \\
1 & 0 & 1 \\
1 & 1 & 0 \\
\bottomrule
\end{tabular}
\end{center}

\medskip
\noindent\textbf{Modelo linear:}
\[
  \hat{y} = x^T w + b = w_1 x_1 + w_2 x_2 + b
\]

\noindent\textbf{Função de perda (MSE):}
\[
  \mathcal{L} = \frac{1}{4}\sum_{i=1}^{4}\left(\hat{y}_i - y_i\right)^2
\]

% ─────────────────────────────────────────────
\section*{2. Formulação Matricial}

Absorvemos o viés $b$ na matriz de design adicionando uma coluna de 1s, de modo que o vetor de parâmetros seja $\theta = [b,\, w_1,\, w_2]^T$:

\[
  X = \begin{bmatrix}
    1 & 0 & 0 \\
    1 & 0 & 1 \\
    1 & 1 & 0 \\
    1 & 1 & 1
  \end{bmatrix},
  \qquad
  y = \begin{bmatrix} 0 \\ 1 \\ 1 \\ 0 \end{bmatrix},
  \qquad
  \theta = \begin{bmatrix} b \\ w_1 \\ w_2 \end{bmatrix}
\]

\noindent A \textbf{equação normal} que minimiza o MSE é:
\[
  X^T X\, \theta = X^T y
\]

% ─────────────────────────────────────────────
\section*{3. Cálculo de $X^T X$}

\[
  X^T = \begin{bmatrix}
    1 & 1 & 1 & 1 \\
    0 & 0 & 1 & 1 \\
    0 & 1 & 0 & 1
  \end{bmatrix}
\]

Cada entrada $(i,j)$ de $X^T X$ é o produto interno da $i$-ésima linha de $X^T$ com a $j$-ésima coluna de $X$:

\[
  X^T X =
  \begin{bmatrix}
    1{+}1{+}1{+}1 & 0{+}0{+}1{+}1 & 0{+}1{+}0{+}1 \\
    0{+}0{+}1{+}1 & 0{+}0{+}1{+}1 & 0{+}0{+}0{+}1 \\
    0{+}1{+}0{+}1 & 0{+}0{+}0{+}1 & 0{+}1{+}0{+}1
  \end{bmatrix}
  =
  \begin{bmatrix}
    4 & 2 & 2 \\
    2 & 2 & 1 \\
    2 & 1 & 2
  \end{bmatrix}
\]

% ─────────────────────────────────────────────
\section*{4. Cálculo de $X^T y$}

\[
  X^T y =
  \begin{bmatrix}
    1{\cdot}0 + 1{\cdot}1 + 1{\cdot}1 + 1{\cdot}0 \\
    0{\cdot}0 + 0{\cdot}1 + 1{\cdot}1 + 1{\cdot}0 \\
    0{\cdot}0 + 1{\cdot}1 + 0{\cdot}1 + 1{\cdot}0
  \end{bmatrix}
  =
  \begin{bmatrix} 2 \\ 1 \\ 1 \end{bmatrix}
\]

% ─────────────────────────────────────────────
\section*{5. Sistema Linear}

Substituindo na equação normal:

\[
  \begin{bmatrix}
    4 & 2 & 2 \\
    2 & 2 & 1 \\
    2 & 1 & 2
  \end{bmatrix}
  \begin{bmatrix} b \\ w_1 \\ w_2 \end{bmatrix}
  =
  \begin{bmatrix} 2 \\ 1 \\ 1 \end{bmatrix}
\]

\noindent Expandindo linha a linha:

\begin{align}
  4b + 2w_1 + 2w_2 &= 2 \tag{I}  \\
  2b + 2w_1 +  w_2 &= 1 \tag{II} \\
  2b +  w_1 + 2w_2 &= 1 \tag{III}
\end{align}

% ─────────────────────────────────────────────
\section*{6. Resolução}

\subsection*{Passo 1 — Subtrair (III) de (II)}

\[
  (\text{II}) - (\text{III}):\quad
  (2w_1 + w_2) - (w_1 + 2w_2) = 0
  \;\Longrightarrow\;
  w_1 - w_2 = 0
  \;\Longrightarrow\;
  \boxed{w_1 = w_2}
\]

\noindent Isso reflete a \emph{simetria do XOR}: as duas features são perfeitamente intercambiáveis.

\subsection*{Passo 2 — Substituir $w_1 = w_2 \equiv w$ em (I)}

\[
  4b + 2w + 2w = 2
  \;\Longrightarrow\;
  4b + 4w = 2
  \;\Longrightarrow\;
  b + w = \tfrac{1}{2}
  \;\Longrightarrow\;
  b = \tfrac{1}{2} - w \tag{IV}
\]

\subsection*{Passo 3 — Substituir $w_1 = w_2 = w$ e (IV) em (II)}

\[
  2\!\left(\tfrac{1}{2} - w\right) + 2w + w = 1
\]
\[
  1 - 2w + 2w + w = 1
\]
\[
  1 + w = 1
\]
\[
  \boxed{w = 0}
\]

\subsection*{Passo 4 — Recuperar $b$}

\[
  b = \tfrac{1}{2} - 0 = \boxed{\tfrac{1}{2}}
\]

% ─────────────────────────────────────────────
\section*{7. Solução Final}

\begin{caixadestaque}
\[
  w_1 = 0, \qquad w_2 = 0, \qquad b = \frac{1}{2}
\]
\[
  \hat{y}(x) = \frac{1}{2} \quad \text{para qualquer entrada } x
\]
\end{caixadestaque}

O modelo colapsa para a previsão constante igual à média dos rótulos:
\[
  \bar{y} = \frac{0 + 1 + 1 + 0}{4} = \frac{1}{2}
\]

O MSE mínimo resultante é:
\[
  \mathcal{L}^* = \frac{1}{4}\left[
    \left(\tfrac{1}{2}\right)^2 +
    \left(\tfrac{1}{2}\right)^2 +
    \left(\tfrac{1}{2}\right)^2 +
    \left(\tfrac{1}{2}\right)^2
  \right] = \frac{1}{4}
\]

% ─────────────────────────────────────────────
\section*{8. Interpretação}

Este resultado não é um acidente numérico — é a \textbf{prova formal de que o XOR não é linearmente separável}. Qualquer gradiente que $w_1$ ou $w_2$ acumulariam a partir de um par de pontos é exatamente cancelado pelo par simétrico. A solução ótima para um modelo linear ignora completamente as features e prevê a média da distribuição.

O MSE de $1/4$ é equivalente ao de um classificador que adivinha aleatoriamente com probabilidade $1/2$ — a melhor hipótese linear não supera o \emph{baseline} trivial.

A saída desta limitação requer \textbf{não-linearidade}: uma rede neural com pelo menos uma camada oculta (e função de ativação não-linear) é capaz de aprender o XOR exatamente.

\end{document}